\section{Single-Block Equivalent and Analysis Definitions}
\label{sec:methods}
\subsection{Single-block equivalent}
\label{sec:equivalent}
Grid studies represent a data center plant by one supply block that is driven by the total server power \cite{ross2025emt,arifin2026dynamic}.
The paper compares this single-block equivalent with the full model of Section~\ref{sec:electrical}.
The comparison should expose only what the equivalent omits, namely the differences between the blocks and the allocation of the tasks to the halls.
The equivalent is therefore built from the parameters and the operating point of the full model alone.

Let the plant contain supply blocks with ratings $S_n$ and weights $w_n=S_n/\sum_nS_n$.
The equivalent block has the rating $S_E=\sum_nS_n$, which is 144~MVA for the 48 blocks of 3~MVA.
Its server power is the total server power of the plant.
A capacitance, a PI gain or the PLL low-pass cut-off $a$ uses the capacity-weighted arithmetic mean $a_E=\sum_nw_na_n$.
A series resistance or inductance uses the capacity-weighted harmonic mean
\begin{equation}
a_E=\left(\sum_n\frac{w_n}{a_n}\right)^{-1}.
\end{equation}
The arithmetic means apply to $c_d$, $c_v$, $c_p$, $c_e$, $\omega_{\rm lp}$ and the 14 PI gains, and the harmonic means apply to $r_a$, $\ell_a$, $r_v$, $\ell_v$ and $r_p$.
The PI gains are averaged directly and are not recomputed from the averaged elements.
The reference values and $\omega_s$ are equal in all blocks and keep their values.
These rules act on device per-unit values, each referred to the rating of its own block.
All blocks of the plant have the same rating, so every weight equals 1/48.

The equivalent keeps one path from the PCC to its block.
The two main transformers, the 12 feeders and the 48 unit transformers each become one branch with the equal-loss impedance
\begin{equation}
Z_E=\frac{\sum_nP_n^2Z_n}{\left(\sum_nP_n\right)^2},
\end{equation}
where $Z_n$ are the member impedances on the system base and $P_n$ are their initial active flows in the full model.
A feeder flow includes the cooling load behind it, whereas a unit-transformer flow includes only its supply block.
The feeder charging susceptances are summed.
The two medium-voltage buses become one bus, and the 12 hall buses and the 48 low-voltage buses each become one bus.
The equal-loss impedance reproduces the branch losses at the operating point when all members share one voltage magnitude and one power factor.
The remaining offset of the initial PCC power is retained and reported in Section~\ref{sec:time}.

The cooling loads of the 12 halls become one constant-power load at the hall bus of the equivalent.
The two auxiliary loads become one constant-impedance load at its medium-voltage bus.
The residual load at bus~8 is that of the full model.

The input of the equivalent is the total server power $P_{\rm IT}=\sum_hP_{{\rm IT},h}$.
A server power in any hall therefore enters the one block in the same way.
Let $g_h$ be the response of the full model to 1~MW of server power in hall $h$, and let $\bar g$ be the response of the equivalent to the same input, which does not depend on $h$.
The full model responds to task $j$ with $\sum_ha_{hj}g_h$, whereas the equivalent responds with $\bar g$ because $\sum_ha_{hj}=1$.
The difference for task $j$ is
\begin{equation}
e_j=\sum_ha_{hj}\left(\bar g-g_h\right).
\label{eq:task-error}
\end{equation}

A positive control checks the construction.
It gives every hall the nominal design, identical feeders and unit transformers and a server power of 6~MW.
Every merged element then equals the parallel combination of identical members, so the equivalent must reproduce the full model.
Over 51 frequencies of the grid of Section~\ref{sec:small-signal}, the relative difference between the responses of the two models to the hall server powers is at most $1.0\cdot10^{-12}$.
The relative difference is the norm of the response difference over the frequencies and the halls divided by the norm of the full-model response, maximized over the four outputs.

\subsection{Small-signal and modal quantities}
\label{sec:small-signal}
Linearization of $\dot x=f(x,z,u)$ and $0=g(x,z)$ gives a state matrix after removal of the common angle reference and elimination of the nonsingular algebraic block.
With derivatives at the operating point,
\begin{equation}
A=f_x-f_zg_z^{-1}g_x,\qquad B=f_u-f_zg_z^{-1}g_u,\qquad G(s)=C(sI-A)^{-1}B+D.
\end{equation}
The output derivatives include the corresponding elimination of $z$.
The operating point is the equilibrium under the mean server power of every hall over the 60~s window of the realization.
The inputs $u$ are the server powers of the halls, which enter the device equations through the commands $p_{\rm load}$ of their blocks.
The four outputs are the PCC active power, reactive power and voltage and the speed of the synchronous machine at bus~1, expressed as a frequency in Hz.
Column $h$ of $G(s)$ is $g_h$, and the task response is $G(s)B_c$ with the allocation matrix $B_c=[a_{hj}]$.

The identical blocks of a hall give repeated eigenvalues.
Eigenvalues closer than $10^{-6}(1+|\lambda|)$ are treated as one repeated eigenvalue.
Repeated modes therefore use an invariant-subspace projector $\Pi_c=V_c(W_c^{\mathsf H}V_c)^{-1}W_c^{\mathsf H}$.
For a repeated eigenvalue $\lambda_c$, the residue of its invariant subspace is $C\Pi_cB$.
Its task-to-output residue follows after multiplication by $B_c$.
The contribution of a mode to a task is the magnitude of its active-power residue divided by $|\operatorname{Re}\lambda_c|$.
A mode counts once regardless of its multiplicity.
An oscillatory mode has a positive imaginary part and a damping ratio below 0.9.
The mode counts of Section~\ref{sec:results} and of Fig.~3(a) of the paper include the oscillatory modes between 3 and 20~Hz whose largest contribution over the tasks reaches 10\% of the largest such contribution of the full model.
In Fig.~3(a) of the paper, the marker area grows with this largest contribution.

The band error measures how closely the equivalent follows the full model in a frequency band.
For output $m$, task $j$ and band $\mathcal B$, it is
\begin{equation}
\epsilon_{mj,\mathcal B}=\left[\frac{\sum_{k\in\mathcal B}|G_{E,mj}(\ii\omega_k)-G_{F,mj}(\ii\omega_k)|^2}{\sum_{k\in\mathcal B}|G_{F,mj}(\ii\omega_k)|^2}\right]^{1/2},
\end{equation}
where $G_{F,mj}$ and $G_{E,mj}$ are the task responses of the full model and of the equivalent.
The grid has 601 logarithmically spaced frequencies from 0.1 to 20~Hz, and the sums weight its points equally.
The low band contains the grid points from 0.1 to 3~Hz, and the resonance band contains the points above 3~Hz up to 20~Hz.
The largest band error of a realization is the maximum over its training and fine-tuning jobs.

Two limiting allocations show how the band error depends on the allocation.
The single-hall band error of hall $h$ uses $g_h$ and $\bar g$ in place of the task responses, which describes a unit input in hall $h$ alone.
The uniform-input band error uses $\frac1H\sum_hg_h$ and $\bar g$, which describes the same unit input spread evenly over the $H=12$ halls.
Both use the PCC active power and the resonance band.

The peak gain and the peak frequency of each model come from a bounded one-dimensional search between the two grid neighbors of the largest grid value of its active-power gain.
The phase difference is $\arg[G_{E,Pj}(\ii\omega_p)/G_{F,Pj}(\ii\omega_p)]$, evaluated at the grid frequency $\omega_p$ of the largest gain of the full model.
The band error also contains phase differences, so a close peak gain can coexist with a large band error.

\subsection{Local supply-block modes and network coupling}
The origin of a resonant response must be distinguished from its transmission to the PCC.
The supply-block equations reproduced from \cite{zhao2026dynamic} use unlimited modulation, a fixed VSI reference frequency $\omega_s$ and a prescribed server power.
Let $\bm v_c^\star$ denote the numerator of \eqref{eq:vsi-modulation}.
Substitution into the VSI current equation \eqref{eq:vsi-current} gives
\begin{equation}
v_d\bm m_v=\bm v_c^\star,\qquad v_d\ne0.
\end{equation}
The command $\bm v_c^\star$ depends only on local VSI states and fixed references.
The DC-link voltage therefore cancels from the VSI state equations.
The eight VSI states, two PSU states and two server DC--DC states form a local 12-state subsystem driven by its own command $p_{\rm load}$.
Network voltage and AFE states do not enter these equations.

This dependence determines how the local modes enter the plant model.
For $N$ supply blocks, let $x_L$ collect their local states and let $x_N$ collect the remaining AFE, DC-link and external-grid states.
After elimination of the network algebraic variables, ordering the states as $x_L,x_N$ gives
\begin{equation}
A=\begin{bmatrix}A_L&0\\A_{NL}&A_N\end{bmatrix},\qquad
A_L=\operatorname{diag}(A_{L,1},\ldots,A_{L,N}),
\end{equation}
where each $A_{L,n}$ is the 12-state matrix of one supply block.
For fixed local operating points and parameters, its eigenvalues do not depend on $A_N$.
Local state disturbances can drive the AFE and DC-link states through $A_{NL}$ and appear at the PCC.
Network coupling remains in $A_N$ and affects the upstream responses.
Thus, a local VSI--PSU mode can appear in several grid-connected quantities without being generated by feedback between supply blocks.
Signed sums of spectral-projector entries are not energy fractions.

The comparisons of the paper and of Section~\ref{sec:results} assess task-to-PCC accuracy.
They do not include a controlled removal or variation of network feedback.
They therefore do not isolate inter-block network coupling as the cause of the dominant resonances or of the differences between the full model and the single-block equivalent.
This structural interpretation applies to the stated unlimited, fixed-frequency model.

\subsection{Time integration and internal variables}
\label{sec:time-methods}
Every model receives the same stored input realization.
A one-second rest precedes the 60-second task window.
Each run starts from the equilibrium of its model at the server power of the first input instant.
Every phase transition appears in the scheduled-event calendar.
The integrator stops at the event, installs the next continuous segment, solves consistent algebraic initial values and resumes.
A finite-width ramp is not substituted for a phase transition.

The production calculation uses Solverz Rodas with KLU, just-in-time compiled modules, relative tolerance $10^{-4}$, absolute tolerance $10^{-6}$, maximum step 0.01~s and output spacing 0.0025~s.
Consistent initialization uses tolerance $10^{-6}$.
Model checks require equal equation and variable counts and initial residual below $10^{-10}$.
The strict comparison repeats both models at relative tolerance $10^{-8}$ and absolute tolerance $10^{-10}$.

For output $y$, set $d_F(t)=y_F(t)-y_F(0)$ and $d_E(t)=y_E(t)-y_E(0)$.
Over the task window, the dynamic error is $\operatorname{RMS}(d_E-d_F)/\operatorname{RMS}(d_F)$.
The initial offset $b_0=y_E(0)-y_F(0)$ is reported separately.
The spectral share of the active-power error comes from a periodogram of $d_E-d_F$ over the task window with the mean removed and one Hann window.
The share between 3 and 9~Hz covers the frequencies above 3~Hz up to 9~Hz.

The internal-variable comparison repeats the production runs of both models under the three inputs of the realization used in the paper.
It records the first supply block of every hall of the full model and the one block of the equivalent.
The four blocks of a hall share their input, parameters and initial state, so the first block represents the hall.
The recorded block variables are $v_d$, the magnitude $|\bm v_v|$ of the VSI output voltage, $v_p$, $v_e$ and the magnitude $|\bm i_a|$ of the AFE current, in per unit on the block base.
The server power $P_{{\rm IT},h}$ and the AC input power $P_{{\rm AC},h}$ of hall $h$ are sums over its four blocks.
The peak deviation of a variable is its largest absolute change from the initial value over the task window.
The ratio divides the peak deviation of the equivalent by the largest peak deviation among the halls of the full model.
For the equivalent, the server power and the AC input power are plant totals, so these two quantities have no ratio.
